\section{Marco teórico}

\subsection{Café (\textit{Coffea})}

\subsubsection{Clasificación botánica del género \textit{Coffea}}

El género \textit{Coffea}, perteneciente a la familia \textit{Rubiaceae} y al orden \textit{Gentianales}, constituye un grupo taxonómico de gran relevancia económica y agroindustrial que comprende al menos 130 especies descritas \parencite{baschenis_comparative_2026, garcia_agro-morphological_2025, park_comparative_2025}. Estos especímenes son botánicamente clasificados como arbustos o pequeños árboles perennes, cuyo origen se localiza primordialmente en las regiones tropicales y subtropicales de África y Asia \parencite{gomez-godinez_exploring_2024, himri_comparative_2025, lorencone_agricultural_2023}. La diversificación del género se ha manifestado de manera independiente en tres centros principales: África Occidental y Central, África Oriental y las islas del Océano Índico occidental, colonizando hábitats que oscilan desde bosques siempreverdes permanentemente húmedos hasta sabanas arboladas estacionalmente secas \parencite{baschenis_comparative_2026}.

Morfológicamente, las plantas de este género se caracterizan por una estructura compleja donde el fruto, una baya denominada habitualmente ``cereza'', está compuesto por el epicarpo (piel externa), el mesocarpo (pulpa y mucílago) y el endocarpo (pergamino), el cual rodea a la semilla \parencite{ferreira_impact_2024, hasan_comparative_2024, osorio_hernandez_bioclimatic_2018}. Las semillas son de naturaleza albuminosa y consisten en un embrión pequeño embebido en un abundante endospermo celular que acumula galactomananos, aceites y metabolitos especializados como cafeína, trigonelina y ácidos clorogénicos \parencite{baschenis_comparative_2026}. La fenología de la planta permite que algunas especies florezcan más de una vez en una misma temporada, produciendo frutos que cambian de color hacia tonos rojos, púrpuras o negros al alcanzar la madurez fisiológica \parencite{osorio_hernandez_bioclimatic_2018, barbosa_sensory_2019, baschenis_comparative_2026}.

\subsubsection{Especies de importancia comercial: arábica y robusta}

En el ámbito comercial, la producción global está dominada por dos especies principales: \textit{Coffea arabica} y \textit{Coffea canephora} (conocida como Robusta), las cuales presentan marcadas diferencias botánicas y agronómicas \parencite{himri_comparative_2025, tieghi_effects_2024}. \textit{Coffea arabica} es una especie alotetraploide que se originó a partir de un evento de hibridación natural entre \textit{Coffea canephora} y \textit{Coffea eugenioides}; representa aproximadamente el 70\% de la producción mundial y es altamente valorada por su calidad sensorial superior, acidez y aroma distintivos \parencite{park_comparative_2025, baschenis_comparative_2026, escuela_superior_politecnica_agropecuaria_de_manabi_manuel_felix_lopez_carrera_de_ingenieria_agricola_postharvest_2023, gallardo-ignacio_chemical_2023}. No obstante, es una planta sensible al incremento de temperaturas, prosperando óptimamente en climas frescos entre los 18~°C y 23~°C, y es vulnerable a patógenos como la roya del café (\textit{Hemileia vastatrix}) \parencite{park_comparative_2025, febrianto_coffee_2023}.

Por el contrario, \textit{Coffea canephora} se distingue por su notable resiliencia ante condiciones ambientales adversas, mostrando una mayor tolerancia al calor y a diversas plagas y enfermedades \parencite{escuela_superior_politecnica_agropecuaria_de_manabi_manuel_felix_lopez_carrera_de_ingenieria_agricola_postharvest_2023, suryono_physicochemical_2024}. Agronómicamente, esta especie es más económica de cultivar debido a su adaptabilidad a climas cálidos, con rangos óptimos de temperatura entre 24~°C y 30~°C \parencite{park_comparative_2025, suryono_physicochemical_2024}. Aunque su perfil sensorial suele considerarse menos complejo que el de \textit{Coffea arabica}, caracterizándose por un sabor más fuerte y mayor cuerpo, posee niveles significativamente más elevados de cafeína, sólidos solubles y compuestos antioxidantes, lo que la hace indispensable para la industria de cafés solubles y mezclas comerciales \parencite{hasan_comparative_2024, escuela_superior_politecnica_agropecuaria_de_manabi_manuel_felix_lopez_carrera_de_ingenieria_agricola_postharvest_2023, himri_comparative_2025, tieghi_effects_2024}.

\subsubsection{Variedades cultivadas en Colombia}

Colombia, tercer productor mundial de café y principal exponente global de 
\textit{Coffea arabica} en zonas de alta montaña, sustenta cerca de 540000 familias 
cafeteras y aporta aproximadamente el 15\% del PIB agropecuario 
\parencite{betancourt_agricultural_2024, gomez-godinez_exploring_2024}. El enfoque 
estratégico hacia \textit{C. arabica} responde a la demanda internacional por su 
calidad sensorial, aroma y acidez característica 
\parencite{escuela_superior_politecnica_agropecuaria_de_manabi_manuel_felix_lopez_carrera_de_ingenieria_agricola_postharvest_2023, 
park_comparative_2025}.

El panorama varietal integra cultivares tradicionales como \textit{Typica}, 
\textit{Bourbon}, \textit{Caturra} y \textit{Mundo Novo}, junto con variedades 
mejoradas por Cenicafé —\textit{Castillo}, \textit{Colombia} y \textit{Costa Rica 
95}— orientadas principalmente a la resistencia contra la roya (\textit{Hemileia 
vastatrix}) \parencite{gallardo-ignacio_chemical_2023, rodriguez_effect_2020, 
febrianto_coffee_2023}. La variedad \textit{Castillo} destaca por su tolerancia a 
la roya en altitudes inferiores a 1600~m.s.n.m. y su elevado potencial productivo 
en zonas de mayor elevación \parencite{rodriguez_effect_2020}. Tanto los materiales 
tradicionales como los mejorados son capaces de producir perfiles sensoriales de 
excelencia —con notas afrutadas, acidez fresca y matices de chocolate y caramelo—, 
cuya expresión depende de la interacción entre la base genética, la variabilidad 
climática y el manejo técnico \parencite{gallardo-ignacio_chemical_2023, 
rodriguez_effect_2020, osorio_hernandez_simulation_2015, tieghi_effects_2024}. 
Así, la diversidad varietal colombiana asegura la resiliencia fitosanitaria de los 
cafetales y garantiza una oferta diferenciada que cumple con los estándares 
internacionales de calidad \parencite{rodriguez_effect_2020, tieghi_effects_2024}.

\subsubsection{Estructura anatómica del fruto y del grano de café}

El fruto del café, clasificado botánicamente como una baya o drupa, está compuesto 
por capas concéntricas que envuelven la semilla \parencite{de_bruyn_exploring_2017, 
jorge_a_d_jesus_molecular_2025}. La capa más externa es el epicarpo, una piel cerosa 
que adquiere tonalidades rojas o amarillas al alcanzar la madurez fisiológica 
\parencite{polakova_changes_2022, ferreira_impact_2024}. Bajo esta se encuentra el 
mesocarpo, conformado por la pulpa carnosa y el mucílago, tejido de parénquima 
constituido por agua (85\%), azúcares y polisacáridos complejos, con un espesor de 
0.5 a 2.0~mm \parencite{ferreira_impact_2024, de_bruyn_exploring_2017, 
jorge_a_d_jesus_molecular_2025, polakova_changes_2022}. El endocarpo o pergamino, 
de naturaleza polisacárida, rodea directamente la semilla ejerciendo funciones de 
resistencia mecánica y protección, mientras que la espermodermis o \textit{silver 
skin} la recubre íntimamente hasta el final del desarrollo del fruto 
\parencite{polakova_changes_2022, baschenis_comparative_2026, 
jorge_a_d_jesus_molecular_2025, de_bruyn_exploring_2017, hasan_comparative_2024}. 
En conjunto, estas capas garantizan la estabilidad higrotérmica y biológica de los 
tejidos internos \parencite{osorio_hernandez_bioclimatic_2018, 
osorio_hernandez_computational_2020}.

La semilla está constituida por un endospermo celular, que representa más del 97\% 
de su masa, y un embrión espatulado embebido en su interior 
\parencite{baschenis_comparative_2026}. El endospermo acumula hasta un 53\% de 
polisacáridos insolubles (celulosa y galactomananos), 18\% de lípidos y 12\% de 
proteínas, además de compuestos bioactivos como cafeína, trigonelina y ácidos 
clorogénicos, que cumplen funciones de protección química de la semilla 
\parencite{ferreira_impact_2024, baschenis_comparative_2026, de_bruyn_exploring_2017}. 
El embrión, aunque representa apenas entre el 0.39\% y 2.87\% de la masa total, es 
el componente vital responsable de la germinación 
\parencite{baschenis_comparative_2026}.

\subsubsection{Café de especialidad: concepto, criterios de calidad y diferenciación}

El café de especialidad se define bajo los estándares de la Specialty Coffee Association 
(SCA) como aquel que alcanza 80 puntos o más en una evaluación sensorial estandarizada, 
diferenciándose del café \textit{commodity} por su excelencia cualitativa, ausencia de 
defectos y capacidad de generar experiencias sensoriales complejas de mayor valor 
comercial \parencite{barbosa_sensory_2019, carmo_quality_2020, febrianto_coffee_2023, 
gallardo-ignacio_chemical_2023, tieghi_effects_2024}. Los criterios de clasificación 
abarcan dimensiones sensoriales —aroma, acidez, cuerpo, sabor, dulzura, balance, 
uniformidad y limpieza de taza— y físicas, exigiéndose ausencia total de defectos 
primarios, cuya presencia compromete tanto el perfil en taza como la inocuidad del 
producto \parencite{febrianto_coffee_2023, barbosa_sensory_2019, carmo_quality_2020, 
dias_alixandre_physical_2023, elhalis_role_2021, tieghi_effects_2024, 
escuela_superior_politecnica_agropecuaria_de_manabi_manuel_felix_lopez_carrera_de_ingenieria_agricola_postharvest_2023}. 
La integridad física del grano está directamente vinculada a su composición química, 
precursora de los compuestos volátiles que definen la identidad del producto durante 
la tostión \parencite{escuela_superior_politecnica_agropecuaria_de_manabi_manuel_felix_lopez_carrera_de_ingenieria_agricola_postharvest_2023, 
tieghi_effects_2024}.

La calidad del café de especialidad emerge de la interacción entre factores genéticos, 
ambientales y técnicos. \textit{Coffea arabica} constituye la base predominante de 
este segmento por su potencial para producir perfiles sensoriales complejos, cuya 
expresión depende del \textit{terroir} —altitud superior a 1000~m.s.n.m., 
microclima y características edáficas— y de prácticas de poscosecha controladas, 
como la cosecha selectiva y procesos de beneficio húmedo, seco o fermentaciones 
innovadoras \parencite{dias_alixandre_physical_2023, ferreira_impact_2024, 
park_comparative_2025, carmo_quality_2020, febrianto_coffee_2023, hasan_comparative_2024, 
tieghi_effects_2024, barbosa_sensory_2019, de_bruyn_exploring_2017, 
ferreira_bernardes_pereira_influence_2020, 
escuela_superior_politecnica_agropecuaria_de_manabi_manuel_felix_lopez_carrera_de_ingenieria_agricola_postharvest_2023}. 
Esta diferenciación se consolida mediante la trazabilidad y el acceso a mercados que 
valoran la transparencia y sostenibilidad en la cadena productiva 
\parencite{dias_alixandre_physical_2023, hasan_comparative_2024, tieghi_effects_2024}.

\subsubsection{Contexto económico nacional e internacional del café}

El café es el segundo producto básico más comercializado a nivel mundial, sustentando 
los medios de vida de aproximadamente 125 millones de personas en más de 50 países 
\parencite{himri_comparative_2025, osorio_hernandez_computational_2020, 
lorencone_agricultural_2023, osorio_hernandez_simulation_2015}. La producción global 
está liderada por Brasil (30--33\%), seguido por Vietnam y Colombia; este último, 
principal exponente de \textit{Coffea arabica} en tierras altas, aporta el 15\% del 
PIB agropecuario, genera cerca de 2.5 millones de empleos y sustenta a 540000 
familias cafeteras —la mayoría en propiedades menores a dos hectáreas—, 
convirtiéndose en una herramienta estratégica para la reducción de la pobreza rural 
\parencite{lorencone_agricultural_2023, hasan_comparative_2024, tieghi_effects_2024, 
betancourt_agricultural_2024}.

La cadena de valor experimenta una transformación hacia la diferenciación cualitativa: 
mientras el café \textit{commodity} crece al 2\% anual, el segmento de especialidad 
avanza al 15\%, impulsado por tendencias globales de sostenibilidad y trazabilidad 
\parencite{barbosa_sensory_2019, febrianto_coffee_2023, tieghi_effects_2024}. Esta 
valorización, que puede elevar el precio de venta hasta 41~USD/kg en lotes de 
excelencia, mitiga la vulnerabilidad de los pequeños productores ante la volatilidad 
de precios internacionales \parencite{febrianto_coffee_2023, dias_alixandre_physical_2023, 
hasan_comparative_2024, tieghi_effects_2024}. No obstante, el sector enfrenta retos 
estructurales como la variabilidad climática —que podría reducir el área apta para 
el cultivo en un 60\% hacia 2050— y la necesidad de innovación tecnológica para 
garantizar la resiliencia fitosanitaria frente a patógenos como la roya 
\parencite{hasan_comparative_2024, lorencone_agricultural_2023, park_comparative_2025}. 
La sostenibilidad del sistema depende así de integrar prácticas avanzadas de 
poscosecha con estrategias de mercado que recompensen la calidad 
\parencite{dias_alixandre_physical_2023, betancourt_agricultural_2024}.

\subsection{Sistema de beneficio del café y ubicación del secado}

\subsubsection{Beneficio húmedo y sus etapas}

El sistema de beneficio húmedo constituye el principal método de procesamiento de café en naciones reconocidas por su alta calidad, tales como Colombia, los países de Centroamérica y Hawái \parencite{hasan_comparative_2024, osorio_hernandez_computational_2020}. A diferencia del proceso seco, este sistema tiene como objetivo primordial la transformación del fruto fresco en café pergamino mediante la remoción sistemática de las capas externas —epicarpo y mesocarpo— a través de medios mecánicos e hídricos \parencite{carmo_sensory_2020, polakova_changes_2022}. Esta metodología es valorada en la industria agroindustrial por su capacidad para resaltar atributos sensoriales superiores, como una mayor acidez y perfiles aromáticos complejos, aunque requiere una infraestructura específica y el uso significativo de agua para garantizar la limpieza y homogeneidad del grano final \parencite{hasan_comparative_2024, polakova_changes_2022}.

La secuencia operativa se inicia con la recolección selectiva manual, etapa crítica donde se seleccionan exclusivamente frutos en estado de madurez ``cereza'' para asegurar que el grano exprese su máximo vigor cualitativo \parencite{barbosa_sensory_2019, ferreira_bernardes_pereira_influence_2020}. Posteriormente, los frutos son sometidos al despulpado, una operación mecánica que desprende la piel y la pulpa, dejando el grano recubierto únicamente por el mucílago viscoso y el endocarpo o pergamino \parencite{polakova_changes_2022, barbosa_sensory_2019}. Esta remoción física expone el interior del fruto a la interacción ambiental, lo que marca el inicio de cambios químicos y microbiológicos que deben ser estrictamente controlados para evitar la proliferación de microorganismos indeseables que podrían degradar la calidad del lote \parencite{carmo_sensory_2020, carmo_quality_2020, de_bruyn_exploring_2017}.

Una de las fases más determinantes es la fermentación, proceso biológico natural donde el mucílago remanente es degradado por la acción de enzimas producidas por levaduras y bacterias epífitas \parencite{polakova_changes_2022, febrianto_coffee_2023, ferreira_impact_2024}. Durante este periodo, que suele oscilar entre 6 y 24 horas, ocurren transformaciones bioquímicas complejas donde los azúcares, lípidos y proteínas se convierten en alcoholes, ésteres y ácidos orgánicos, los cuales actúan como precursores del sabor y el aroma \parencite{polakova_changes_2022, ferreira_impact_2024, carmo_sensory_2020}. Una fermentación controlada no solo facilita el desprendimiento del mucílago, sino que permite la síntesis de notas dulces, cítricas y frutales que añaden valor agregado al producto \parencite{carmo_sensory_2020}. Una vez finalizada la degradación biológica, se procede al lavado con agua limpia para remover los residuos de la fermentación e impurezas, deteniendo así las reacciones bioquímicas en el punto óptimo \parencite{hasan_comparative_2024, barbosa_sensory_2019, dias_alixandre_physical_2023}.

Finalmente, el proceso de beneficio húmedo culmina con la etapa de secado, cuyo propósito es reducir el contenido de humedad del grano desde niveles iniciales del 55--60\% hasta un rango de equilibrio de entre el 10\% y 12\% \parencite{barbosa_sensory_2019, de_bruyn_exploring_2017, osorio_hernandez_simulation_2015}. Este paso final es vital para estabilizar el material biológico, inhibir procesos de deterioro y garantizar la inocuidad durante el almacenamiento prolongado, previniendo defectos físicos y sensoriales \parencite{osorio_hernandez_simulation_2015}. La excelencia en la taza depende, por tanto, de un control riguroso en cada una de estas etapas, ya que cualquier error de diseño o manejo en el beneficio puede comprometer la composición química intrínseca del grano y, en última instancia, su clasificación como café de especialidad \parencite{escuela_superior_politecnica_agropecuaria_de_manabi_manuel_felix_lopez_carrera_de_ingenieria_agricola_postharvest_2023, ferreira_bernardes_pereira_influence_2020, osorio_hernandez_computational_2020}.

% =============================================================
% COLORES
% =============================================================
\definecolor{cafeOscuro}{RGB}{101, 67, 33}
\definecolor{cafeClaro}{RGB}{210, 180, 140}
\definecolor{colorNatural}{RGB}{193, 120, 40}
\definecolor{colorHoney}{RGB}{160, 130, 10}
\definecolor{colorLavado}{RGB}{50, 110, 175}
\definecolor{colorComun}{RGB}{70, 125, 75}
\definecolor{colorDecision}{RGB}{130, 60, 150}
\definecolor{bgNatural}{RGB}{255, 243, 215}
\definecolor{bgHoney}{RGB}{255, 250, 210}
\definecolor{bgLavado}{RGB}{220, 238, 255}
\definecolor{bgComun}{RGB}{228, 244, 228}
\definecolor{bgDecision}{RGB}{245, 230, 250}

% =============================================================
% ESTILOS
% =============================================================
\tikzset{
    etapaComun/.style={
        rectangle, rounded corners=5pt,
        draw=colorComun, line width=1.1pt, fill=bgComun,
        text width=3.2cm, minimum height=0.9cm,
        align=center, font=\small\bfseries, text=cafeOscuro
    },
    decision/.style={
        diamond, aspect=2.6,
        draw=colorDecision, line width=1.1pt, fill=bgDecision,
        text width=2.4cm, align=center,
        font=\footnotesize\bfseries, text=colorDecision, inner sep=1pt
    },
    etapaNatural/.style={
        rectangle, rounded corners=5pt,
        draw=colorNatural, line width=1.1pt, fill=bgNatural,
        text width=2.4cm, minimum height=0.85cm,
        align=center, font=\small\bfseries, text=cafeOscuro
    },
    etapaHoney/.style={
        rectangle, rounded corners=5pt,
        draw=colorHoney, line width=1.1pt, fill=bgHoney,
        text width=2.4cm, minimum height=0.85cm,
        align=center, font=\small\bfseries, text=cafeOscuro
    },
    etapaLavado/.style={
        rectangle, rounded corners=5pt,
        draw=colorLavado, line width=1.1pt, fill=bgLavado,
        text width=2.4cm, minimum height=0.85cm,
        align=center, font=\small\bfseries, text=cafeOscuro
    },
    productoFinal/.style={
        rectangle, rounded corners=5pt,
        draw=cafeOscuro, line width=1.4pt, fill=cafeClaro,
        text width=3.2cm, minimum height=0.9cm,
        align=center, font=\small\bfseries, text=cafeOscuro
    },
    flechaComun/.style={
        -{Stealth[length=5pt,width=4pt]},
        line width=1.0pt, color=cafeOscuro
    },
    flechaNatural/.style={
        -{Stealth[length=5pt,width=4pt]},
        line width=1.0pt, color=colorNatural
    },
    flechaHoney/.style={
        -{Stealth[length=5pt,width=4pt]},
        line width=1.0pt, color=colorHoney
    },
    flechaLavado/.style={
        -{Stealth[length=5pt,width=4pt]},
        line width=1.0pt, color=colorLavado
    },
    labelSiNo/.style={
        font=\scriptsize\itshape, fill=white,
        inner sep=1.5pt, text=gray
    },
}

% =============================================================

\begin{figure}[H]
\centering
\caption{Diagrama de flujo del proceso de beneficio del café
         (\textit{Coffea arabica} L.)}
\label{fig:diagrama_beneficio}
\vspace{0.3em}

% Coordenadas absolutas para control total del layout
\begin{tikzpicture}[x=1cm, y=1cm]

    % ----------------------------------------------------------
    % NODOS — posiciones absolutas (xshift desde el centro = 0)
    % Eje central en x = 0
    % Rama izquierda centrada en x = -4.5
    % Rama derecha centrada en x = +4.8
    % ----------------------------------------------------------

    % Eje central
    \node[etapaComun]          at (0,  0.0) (recoleccion)  {Recolección};
    \node[etapaComun]          at (0, -1.7) (clasificacion){Clasificación y limpieza};
    \node[decision]            at (0, -3.5) (dec1)         {\footnotesize¿Se despulpa\\el fruto?};
    \node[etapaComun]          at (0, -5.3) (despulpado)   {Despulpado};
    \node[decision]            at (0, -7.1) (dec2)         {\footnotesize¿Se retira\\el mucílago?};
    \node[etapaComun]          at (0,-11.0) (secado)        {Secado\\{\footnotesize 10\,--\,12\,\% b.h.}};
    \node[productoFinal]       at (0,-12.8) (producto)     {Café pergamino seco};

    % Rama izquierda
    \node[etapaNatural]        at (-4.5, -3.5) (natural)   {Fruto completo};
    \node[etapaHoney]          at (-4.5, -7.1) (honey)     {Mucílago adherido\\al pergamino};

    % Rama derecha
    \node[etapaLavado]         at ( 4.8, -7.1) (fermentacion) {Fermentación};
    \node[etapaLavado]         at ( 4.8, -9.0) (lavadoAgua)   {Lavado con agua};

    % ----------------------------------------------------------
    % FLECHAS EJE CENTRAL
    % ----------------------------------------------------------
    \draw[flechaComun] (recoleccion)   -- (clasificacion);
    \draw[flechaComun] (clasificacion) -- (dec1);
    \draw[flechaComun] (dec1.south)    -- (despulpado.north)
        node[labelSiNo, pos=0.45, right=2pt]{Sí};
    \draw[flechaComun] (despulpado)    -- (dec2);
    \draw[flechaComun] (secado)        -- (producto);

    % ----------------------------------------------------------
    % BIFURCACIÓN 1 — Natural  (sale oeste de dec1)
    % ----------------------------------------------------------
    \draw[flechaNatural]
        (dec1.west) -- (natural.east)
        node[labelSiNo, midway, above=1pt,
             font=\scriptsize\itshape, text=colorNatural]{No};

    % ----------------------------------------------------------
    % BIFURCACIÓN 2 — Honey  (sale oeste de dec2)
    % ----------------------------------------------------------
    \draw[flechaHoney]
        (dec2.west) -- (honey.east)
        node[labelSiNo, midway, above=1pt,
             font=\scriptsize\itshape, text=colorHoney]{No};

    % ----------------------------------------------------------
    % BIFURCACIÓN 2 — Lavado  (sale este de dec2)
    % ----------------------------------------------------------
    \draw[flechaLavado]
        (dec2.east) -- (fermentacion.west)
        node[labelSiNo, midway, above=1pt,
             font=\scriptsize\itshape, text=colorLavado]{Sí};
    \draw[flechaLavado] (fermentacion) -- (lavadoAgua);

    % ----------------------------------------------------------
    % RUTAS HACIA SECADO
    %
    % NATURAL:
    %   natural.west → izquierda hasta x = -6.2
    %   baja hasta y = -11.0 (misma Y que secado)
    %   derecha hasta secado.west  (entra Y +0.18 arriba del centro)
    %
    % HONEY:
    %   honey.west → izquierda hasta x = -5.7  (distinto a -6.2)
    %   baja hasta y = -11.0
    %   derecha hasta secado.west  (entra justo en el centro)
    %
    % LAVADO:
    %   lavadoAgua.south → baja hasta y = -11.0
    %   izquierda hasta secado.east
    % ----------------------------------------------------------

    % ---- NATURAL ----
    % Segmentos: oeste → x=-6.2 | baja → y=-11.0+0.18 | este → secado.west+0.18
    \draw[flechaNatural]
        (natural.west)
        -- (-6.5, -3.5)
        -- (-6.5, -10.82)
        -- ($(secado.west) + (0, 0.18)$);

    % ---- HONEY ----
    % Segmentos: oeste → x=-5.7 | baja → y=-11.0 | este → secado.west
    \draw[flechaHoney]
        (honey.west)
        -- (-6.0, -7.1)
        -- (-6.0, -11.0)
        -- (secado.west);

    % ---- LAVADO ----
    % Segmentos: sur de lavadoAgua → y=-11.0 | oeste → secado.east
    \draw[flechaLavado]
        (lavadoAgua.south)
        -- (4.8, -11.0)
        -- ($(secado.east) + (0.4, 0)$)
        -- (secado.east);

    % ----------------------------------------------------------
    % RECUADROS PUNTEADOS + ETIQUETAS
    % ----------------------------------------------------------
    \begin{scope}[on background layer]
        \node[
            draw=colorNatural, line width=0.8pt,
            fill=bgNatural, fill opacity=0.35,
            dashed, rounded corners=7pt,
            fit=(natural), inner sep=0.4cm
        ] (boxNatural) {};

        \node[
            draw=colorHoney, line width=0.8pt,
            fill=bgHoney, fill opacity=0.35,
            dashed, rounded corners=7pt,
            fit=(honey), inner sep=0.4cm
        ] (boxHoney) {};

        \node[
            draw=colorLavado, line width=0.8pt,
            fill=bgLavado, fill opacity=0.35,
            dashed, rounded corners=7pt,
            fit=(fermentacion)(lavadoAgua), inner sep=0.5cm
        ] (boxLavado) {};
    \end{scope}

    \node[above=0.05cm of boxNatural,
          font=\scriptsize\bfseries, text=colorNatural,
          fill=white, inner sep=1.5pt] {Método Natural};

    \node[above=0.05cm of boxHoney,
          font=\scriptsize\bfseries, text=colorHoney,
          fill=white, inner sep=1.5pt] {Método Honey};

    \node[above=0.05cm of boxLavado,
          font=\scriptsize\bfseries, text=colorLavado,
          fill=white, inner sep=1.5pt] {Método Lavado};

\end{tikzpicture}

\vspace{0.4em}
\small\textit{Fuente: Adaptado de \parencite{meja_investigating_2025, febrianto_coffee_2023, polakova_changes_2022}. Elaboración propia a partir de las referencias citadas.}

\end{figure}

\subsubsection{Flujo del proceso en plantas de beneficio}

El flujo del proceso en una planta de beneficio húmedo comprende una cadena de fases 
interconectadas orientadas a transformar el fruto fresco en café pergamino seco, 
preservando su integridad biológica y potencial sensorial 
\parencite{de_bruyn_exploring_2017, hasan_comparative_2024}. La secuencia inicia con 
la clasificación hidráulica por densidad para separar frutos maduros de impurezas y 
flotantes, seguida del despulpado mecánico que remueve el epicarpo y parte del 
mesocarpo, exponiendo el grano a condiciones ambientales que desencadenan procesos 
bioquímicos determinantes para la calidad final \parencite{carmo_quality_2020, 
dias_alixandre_physical_2023, ferreira_impact_2024, barbosa_sensory_2019, 
carmo_sensory_2020, osorio_hernandez_computational_2020, 
osorio_hernandez_simulation_2015}. Posteriormente, la fermentación espontánea en 
tanques degrada el mucílago remanente, el lavado con agua limpia detiene las 
reacciones enzimáticas en su punto óptimo, y finalmente el secado —en patios, camas 
suspendidas o sistemas mecánicos— reduce la humedad hasta niveles de equilibrio 
estable \parencite{de_bruyn_exploring_2017, dias_alixandre_physical_2023, 
barbosa_sensory_2019, guerra_garcia_bioclimatic_2022}.

La eficiencia del sistema depende de una gestión integral donde la configuración 
de la planta y la disposición de equipos eviten interrupciones que induzcan 
fermentaciones indeseadas \parencite{osorio_hernandez_computational_2020, 
osorio_hernandez_bioclimatic_2018}. Un control riguroso del flujo higrotérmico es 
esencial, ya que cuellos de botella en la transición entre etapas o una deficiente 
gestión de las cargas térmicas de las secadoras pueden generar defectos físicos y 
sensoriales que degradan el valor comercial del producto \parencite{barbosa_sensory_2019, 
osorio_hernandez_bioclimatic_2018, tieghi_effects_2024, guerra_garcia_bioclimatic_2022, 
osorio_hernandez_computational_2020}. Así, la planta de beneficio actúa como un 
entorno de control metabólico donde cada operación debe sincronizarse para preservar 
los precursores químicos de aroma y sabor hasta alcanzar la estabilidad final del 
pergamino \parencite{escuela_superior_politecnica_agropecuaria_de_manabi_manuel_felix_lopez_carrera_de_ingenieria_agricola_postharvest_2023, 
osorio_hernandez_simulation_2015, tieghi_effects_2024}.

\subsubsection{El secado como punto crítico de control}

El secado del café se constituye como la operación final y determinante del sistema de beneficio húmedo, cuya función primordial es la estabilización biológica y química del grano mediante la reducción controlada de su contenido de humedad desde niveles del 55--60\% hasta un rango de equilibrio de entre el 10\% y 12\% como se mencionó anteriormente \parencite{osorio_hernandez_simulation_2015}. Esta etapa es fundamental para preservar las características cualitativas del producto y permitir su almacenamiento prolongado, ya que la disminución de la actividad de agua inhibe los procesos metabólicos celulares y el deterioro prematuro del endospermo \parencite{osorio_hernandez_simulation_2015, osorio_hernandez_bioclimatic_2018}. Desde una perspectiva técnica, el secado no solo busca la remoción de agua, sino que actúa como una fase de consolidación de los precursores químicos de sabor y aroma sintetizados en etapas previas, garantizando que el potencial genético y sensorial de la variedad sea transferido íntegramente a la bebida final \parencite{barbosa_sensory_2019, de_bruyn_exploring_2017}.

Dada su influencia directa sobre la integridad física y organoléptica, el secado es clasificado como un punto crítico de control (PCC) dentro de la cadena de procesamiento agroindustrial. Un manejo inadecuado en esta fase conlleva riesgos severos, tales como la proliferación de microorganismos patógenos, destacando hongos filamentosos como \textit{Aspergillus ochraceus}, el cual es responsable de la producción de micotoxinas (específicamente la ocratoxina A) que comprometen la inocuidad y seguridad alimentaria del producto \parencite{osorio_hernandez_bioclimatic_2018, osorio_hernandez_computational_2020}. Asimismo, un secado deficiente puede inducir la degradación de compuestos de calidad mediante la peroxidación de lípidos y la desestructuración de las membranas celulares, lo que genera defectos físicos como el blanqueamiento o decoloración del grano y atributos sensoriales negativos en taza, tales como sabores fermentados, sucios o amargos \parencite{escuela_superior_politecnica_agropecuaria_de_manabi_manuel_felix_lopez_carrera_de_ingenieria_agricola_postharvest_2023, osorio_hernandez_simulation_2015, taveira_perfis_2012}.

La eficacia de este proceso depende del control riguroso de variables operativas como la temperatura del aire, el tiempo de exposición y la humedad relativa del ambiente de procesamiento \parencite{ferreira_bernardes_pereira_influence_2020, taveira_perfis_2012}. Para el café de especialidad, es imperativo que la temperatura no exceda rangos críticos (generalmente situados entre 39~°C y 40~°C para café pergamino), ya que niveles superiores pueden provocar la muerte térmica del embrión y acelerar la descomposición de proteínas y carbohidratos esenciales \parencite{carmo_quality_2020, osorio_hernandez_simulation_2015, taveira_perfis_2012}. La gestión técnica debe integrar estrategias de ventilación eficientes para evacuar el vapor de agua generado, evitando la saturación higrotérmica dentro de las instalaciones de beneficio, lo cual es vital para prevenir la rehumectación del grano y asegurar la estabilidad de los cafés con altos puntajes de excelencia \parencite{guerra_garcia_bioclimatic_2022, osorio_hernandez_computational_2020, osorio_hernandez_simulation_2015}.

\subsubsection{Limitaciones tecnológicas actuales en plantas de beneficio}

Las plantas de beneficio de café tradicionales y semi-tecnificadas presentan limitaciones tecnológicas estructurales que comprometen la estandarización y excelencia del producto. Uno de los problemas más críticos es la variabilidad en el control de procesos biológicos como la fermentación, la cual suele realizarse de forma espontánea y sin monitoreo técnico riguroso en finca, lo que deriva en una falta de predictibilidad sobre la calidad final de la taza \parencite{carmo_sensory_2020, polakova_changes_2022}. Esta inconsistencia se ve agravada por una alta dependencia de las condiciones ambientales; dado que la cosecha en países como Colombia suele coincidir con temporadas de lluvia, la humedad relativa externa superior al 70\% constituye una limitación bioclimática natural para las instalaciones que dependen de ventilación pasiva \parencite{osorio_hernandez_bioclimatic_2018, osorio_hernandez_simulation_2015}. La incapacidad de regular variables fundamentales como la temperatura, el tiempo y la humedad durante el procesamiento e incubación induce riesgos de deterioro físico, tales como el blanqueamiento del grano y la degradación de precursores químicos de aroma y sabor \parencite{osorio_hernandez_simulation_2015}.

Desde la perspectiva de la eficiencia de recursos y el impacto ambiental, el sistema de beneficio húmedo convencional exhibe un elevado consumo hídrico, llegando a demandar volúmenes significativos de agua para la remoción del mucílago, lo que genera efluentes con alta carga orgánica y contaminantes difíciles de gestionar sin tecnologías de tratamiento adecuadas \parencite{polakova_changes_2022, escuela_superior_politecnica_agropecuaria_de_manabi_manuel_felix_lopez_carrera_de_ingenieria_agricola_postharvest_2023, rodriguez_effect_2020}. En el ámbito energético, la etapa de secado representa el mayor punto de consumo de potencia en la postcosecha, operando habitualmente con eficiencias térmicas de apenas el 50\% en sistemas mecánicos \parencite{osorio_hernandez_bioclimatic_2018, osorio_hernandez_computational_2020}. Esta ineficiencia provoca la acumulación de calor sensible y vapor de agua dentro de la infraestructura de beneficio, elevando la temperatura interna y la humedad relativa por encima de los límites de estabilidad biológica recomendados (20~°C y 60--70\% HR). Como consecuencia, se incrementa el riesgo de proliferación de microorganismos patógenos, como el hongo \textit{Aspergillus ochraceus}, y se compromete la estabilidad física y sensorial del café pergamino durante su almacenamiento \parencite{osorio_hernandez_bioclimatic_2018, osorio_hernandez_simulation_2015}.

La persistencia de estos métodos artesanales y la baja tasa de tecnificación dificultan la transición hacia el segmento de cafés de especialidad, el cual exige una ausencia total de defectos y perfiles sensoriales altamente reproducibles \parencite{barbosa_sensory_2019, febrianto_coffee_2023}. El manejo actual es predominantemente manual e intensivo en mano de obra, lo que limita la capacidad del productor para competir en mercados que valoran la trazabilidad y la consistencia técnica \parencite{carmo_quality_2020, febrianto_coffee_2023}. 

Existe, por tanto, una necesidad crítica de modernización agroindustrial que trascienda el modelo de \textit{commodity}. La incorporación de innovaciones controladas y sistemas de secado con monitoreo de variables fisicoquímicas es esencial no solo para optimizar el uso de energía y agua, sino para asegurar la preservación del potencial genético de la variedad y la sostenibilidad económica del sistema productivo frente a la variabilidad climática global \parencite{febrianto_coffee_2023, hasan_comparative_2024}.

\subsubsection{Transferencia simultánea de calor y masa}

El secado es una operación unitaria que reduce el contenido de humedad del material 
biológico hasta niveles de equilibrio higroscópico, inhibiendo el crecimiento 
microbiano y las reacciones enzimáticas degradativas 
\parencite{abbaspour-gilandeh_effect_2021, abidin_improving_2024, 
coelho_physicochemical_2024, cruz-ospina_improving_2026}. Físicamente, consiste en 
la transferencia simultánea y acoplada de calor y masa entre el grano y el medio 
circundante: la energía térmica suministrada eleva la temperatura del producto y 
aporta el calor latente necesario para romper las fuerzas de unión del agua en la 
estructura celular, permitiendo su transición de fase líquida a vapor 
\parencite{fauzi_experimental_2025, geng_thermodynamically_2023, yu_current_2025, 
abidin_improving_2024, li_study_2021}. La transferencia de calor ocurre por 
convección —mecanismo predominante en sistemas convectivos—, conducción hacia el 
interior del grano y radiación solar o infrarroja, siendo el gradiente térmico entre 
el aire y el material la fuerza impulsora que determina el potencial de evaporación 
\parencite{chojnacka_improvements_2021, yu_current_2025, nafisah_implementation_2024, 
liu_optimization_2025, abdissa_harmonizing_2023, cruz-ospina_improving_2026, 
institute_of_food_and_biotechnology_can_tho_university_vietnam_impact_2022, 
p_romero_leiton_mathematical_2025}.

La transferencia de masa describe el movimiento de la humedad desde el núcleo del 
grano hacia su superficie, impulsado por gradientes de concentración y diferencias 
de presión de vapor, mediante difusión líquida y acción capilar 
\parencite{yu_current_2025, cruz-ospina_improving_2026, hadibi_effect_2025, 
fauzi_experimental_2025}. Este transporte enfrenta dos tipos de resistencias: la 
externa, asociada a la capa límite de aire que rodea al producto, y la interna, 
propia de la estructura celular del café, que se intensifica conforme avanza el 
secado por fenómenos como la contracción, la reducción de porosidad y el 
endurecimiento superficial (\textit{case hardening}), limitando la migración del 
agua ligada y reduciendo la tasa de secado en las etapas finales 
\parencite{bustos-vanegas_numerical_2022, tejeda-miramontes_phase-segmented_2025, 
geng_thermodynamically_2023, liu_optimization_2025, abbaspour-gilandeh_effect_2021, 
abreu_influence_2025}.

\subsubsection{Mecanismos de migración de humedad}

En el grano de café, el agua se distribuye en tres categorías: libre, ligada y 
citoplasmática. El agua libre, ubicada en la superficie y macroporos, se evapora 
fácilmente debido a su débil unión con la matriz sólida; el agua ligada, atrapada 
en microestructuras celulares, requiere mayor energía para su migración 
\parencite{abidin_improving_2024, alves_physiological_2017}. El transporte interno 
de humedad es dinámico y varía según el avance del secado: en las etapas iniciales 
predomina el flujo capilar, que desplaza el agua líquida hacia la superficie por 
tensión superficial; al agotarse el agua libre, el proceso entra en el periodo de 
velocidad decreciente, donde la difusión líquida —por gradientes de concentración 
en la matriz sólida— y la difusión de vapor —por transición de fase en los espacios 
porosos— se convierten en los mecanismos limitantes 
\parencite{kilic_preservation_2024, yu_current_2025, fernando_effect_2025, 
geng_thermodynamically_2023, hadibi_effect_2025}.

Conforme avanza la deshidratación, la estructura celular experimenta transformaciones 
críticas: la contracción (\textit{shrinkage}) aumenta la densidad y tortuosidad de 
los caminos de transporte, dificultando la extracción del agua residual, mientras 
que el endurecimiento superficial (\textit{case hardening}) genera una capa compacta 
de baja permeabilidad que incrementa drásticamente la resistencia al transporte de 
masa \parencite{kilic_preservation_2024, abidin_improving_2024, 
tejeda-miramontes_phase-segmented_2025, liu_optimization_2025}. Estos cambios 
morfológicos, sumados a la reducción de los gradientes de concentración hacia el 
final del proceso, explican por qué la eliminación de agua se vuelve progresivamente 
más lenta hasta alcanzar el equilibrio higroscópico 
\parencite{amaral_simulation_2018, yu_current_2025}.

\subsubsection{Variables operativas del proceso}

El secado de café está gobernado por variables operativas interdependientes que 
determinan la eficiencia energética y la integridad del producto. La temperatura 
del aire es el parámetro más influyente: incrementa el gradiente térmico, reduce 
la viscosidad del agua y acelera su difusión a través de la estructura porosa del 
grano \parencite{abbaspour-gilandeh_effect_2021, 
institute_of_food_and_biotechnology_can_tho_university_vietnam_impact_2022, 
alves_influence_2020, phitakwinai_thinlayer_2019}. Este efecto está acoplado a la 
humedad relativa (HR): valores bajos de HR intensifican el gradiente de presión de 
vapor entre la superficie del material y el aire, aumentando la fuerza impulsora de 
evaporación; a su vez, temperaturas más altas tienden a reducir la HR, optimizando 
el potencial de secado del sistema \parencite{alves_physiological_2017, 
phitakwinai_thinlayer_2019, gomez-daza_innovacion_2023, p_romero_leiton_mathematical_2025}. 
La velocidad del aire mejora la transferencia convectiva —especialmente en las etapas 
iniciales—, mientras que el espesor de la capa condiciona la uniformidad del proceso: 
capas delgadas reducen la resistencia interna y el tiempo de secado, en tanto que 
capas excesivamente gruesas generan secado heterogéneo y favorecen procesos 
fermentativos indeseables por retención prolongada de humedad 
\parencite{alves_physiological_2017, amaral_simulation_2018, abdissa_harmonizing_2023}.

La gestión inadecuada de estas variables compromete directamente la calidad del 
producto. Superar umbrales críticos de temperatura (40--45~°C) causa daños térmicos 
irreversibles como degradación de compuestos volátiles, oscurecimiento del grano y 
pérdida de capacidad antioxidante; un secado excesivamente prolongado favorece el 
crecimiento microbiano, mientras que el \textit{over-drying} genera granos 
quebradizos y pérdidas económicas por reducción de peso 
\parencite{abdissa_harmonizing_2023, alves_physiological_2017, coelho_physicochemical_2024, 
abidin_improving_2024, abreu_influence_2025}. Por ello, la implementación de sistemas 
de control como controladores PID o lógica difusa es esencial para ajustar 
dinámicamente las condiciones del aire en tiempo real, garantizando el equilibrio 
higroscópico óptimo y preservando los atributos de especialidad del café 
\parencite{arifin_controlling_2024, nafisah_implementation_2024, yu_current_2025, 
coelho_physicochemical_2024, ordonez-lozano_monitoring_2025}.

\subsubsection{Sistema de control en diversos tipos de secado}

El control de variables operativas —temperatura, humedad relativa y flujo de aire— 
es determinante para la tasa de remoción de agua y la integridad del producto 
\parencite{abreu_influence_2025, arifin_controlling_2024}. El secado solar depende 
de un control manual o empírico, generando procesos lentos, heterogéneos y 
susceptibles a contaminación; el secado mecánico convencional introdujo termostatos 
y control ON/OFF, aunque estos esquemas presentan oscilaciones significativas 
respecto al \textit{set-point} y carecen de respuesta ante cambios sutiles en la 
dinámica del proceso \parencite{abreu_influence_2025, arifin_controlling_2024, 
p_romero_leiton_mathematical_2025, bustos-vanegas_numerical_2022, 
cruz-ospina_improving_2026}. Para superar estas limitaciones, se han implementado 
estrategias avanzadas como el control PID —que minimiza el error y estabiliza el 
sistema frente a perturbaciones—, la lógica difusa —que gestiona relaciones no 
lineales mediante reglas lingüísticas— y redes neuronales artificiales, que permiten 
ajustes dinámicos en tiempo real optimizando variables como la potencia infrarroja 
o la velocidad de ventiladores \parencite{arifin_controlling_2024, yu_current_2025, 
gradov_modelling_2022, nafisah_implementation_2024, el-mesery_application_2025-1}.

En sistemas avanzados como los secadores asistidos por bomba de calor, el control 
trasciende la regulación de temperatura para manipular el par 
temperatura--humedad relativa, definiendo con precisión el potencial de secado y 
garantizando una remoción de humedad homogénea \parencite{amaral_simulation_2018, 
bustos-vanegas_numerical_2022, gomez-daza_innovacion_2023}. Esta gestión integral 
reduce significativamente el consumo energético específico (SEC) y el tiempo de 
procesamiento, preservando compuestos volátiles esenciales para el valor comercial 
del producto \parencite{abbaspour-gilandeh_effect_2021, arifin_controlling_2024, 
cruz-ospina_improving_2026}. La incorporación de monitoreo mediante el Internet de 
las Cosas (IoT) asegura la consistencia entre lotes y minimiza el impacto ambiental 
del proceso \parencite{yu_current_2025, p_romero_leiton_mathematical_2025}.

\subsubsection{Control por humedad relativa como variable de proceso}

La humedad relativa (HR) actúa como la variable determinante en la transferencia de masa durante el secado, al definir el grado de saturación del aire y, por ende, su capacidad intrínseca para absorber y transportar el agua proveniente del material \parencite{abidin_improving_2024, gomez-daza_innovacion_2023}. Físicamente, la HR controla el gradiente de presión de vapor entre la superficie del producto y el aire circundante; una reducción en la HR disminuye la presión parcial de vapor en el ambiente, intensificando la fuerza impulsora que acelera la evaporación de la humedad contenida en los tejidos \parencite{cruz-ospina_improving_2026, phitakwinai_thinlayer_2019}. Este mecanismo establece el límite físico del proceso, ya que la remoción de agua cesa cuando el material alcanza el equilibrio higroscópico, estado en el cual la presión de vapor en la superficie del grano se iguala a la del aire, impidiendo cualquier intercambio adicional de masa \parencite{phitakwinai_thinlayer_2019, cruz-ospina_improving_2026}. Por consiguiente, la gestión del par temperatura--humedad relativa resulta más efectiva que el control exclusivo de la temperatura, ya que permite manipular directamente el potencial de secado y mantener una fuerza impulsora constante, independientemente de las fluctuaciones térmicas externas \parencite{gomez-daza_innovacion_2023}.

La regulación precisa de la humedad relativa es fundamental para mitigar fenómenos indeseables que comprometen la integridad fisicoquímica y el valor comercial del producto. Un control inadecuado, caracterizado por niveles de HR excesivamente bajos y temperaturas altas, puede inducir el endurecimiento superficial (\textit{case hardening}), donde la rápida evaporación externa provoca cambios estructurales y consolidación de la matriz que atrapan el agua residual en el núcleo del material, reduciendo drásticamente su permeabilidad y difusividad \parencite{liu_optimization_2025, tejeda-miramontes_phase-segmented_2025}. Asimismo, el monitoreo constante de la HR previene el sobresecado, que degrada compuestos volátiles y aumenta la fragilidad del grano, y evita la rehumectación del producto causada por incrementos súbitos en la humedad del aire de entrada \parencite{abdissa_harmonizing_2023, hadibi_effect_2025, p_romero_leiton_mathematical_2025}. En los sistemas modernos, como los secadores asistidos por bomba de calor (SBC) o sistemas híbridos, esta precisión se alcanza mediante la deshumidificación activa del aire —usando evaporadores o materiales desecantes—, la recirculación controlada del flujo y el uso de circuitos cerrados \parencite{chojnacka_improvements_2021, gomez-daza_innovacion_2023, rahnama_development_2025}. Estas tecnologías permiten operar bajo condiciones donde el fenómeno de secado queda gobernado predominantemente por la dinámica difusional asociada a la humedad relativa, garantizando una remoción de agua homogénea, eficiente y orientada a la preservación de la calidad sensorial \parencite{abdissa_harmonizing_2023, gomez-daza_innovacion_2023}.

\subsection{Cinética de secado y modelamiento matemático}

\subsubsection{Curvas de secado}

La curva de secado describe la evolución del contenido de humedad o de la velocidad de remoción de agua en función del tiempo \parencite{hadibi_effect_2025, yu_current_2025}. En materiales biológicos, se emplean principalmente curvas de razón de humedad (MR) y de velocidad de secado, las cuales reflejan la transferencia simultánea de calor y masa durante el proceso \parencite{guemouni_convective_2022, geng_thermodynamically_2023}. En el caso del café, el secado ocurre predominantemente en el periodo de velocidad decreciente, donde el fenómeno está gobernado por la difusión interna y la resistencia del tejido a la salida del agua \parencite{fernando_effect_2025, alves_physiological_2017}.

El análisis de estas curvas permite evaluar el efecto de variables operativas como la temperatura, la humedad relativa y la velocidad del aire \parencite{abdissa_harmonizing_2023, phitakwinai_thinlayer_2019}. En particular, el incremento de la temperatura acelera la cinética de secado al aumentar la difusividad efectiva y reducir el tiempo de proceso, mientras que factores como el espesor de la capa y la humedad del aire influyen en la fuerza impulsora de evaporación \parencite{fauzi_influence_2024, phitakwinai_thinlayer_2019}. El proceso tiende a estabilizarse al alcanzar la humedad de equilibrio, donde cesa la transferencia neta de masa \parencite{cruz-ospina_improving_2026}.

Desde el punto de vista ingenieril, estas curvas son fundamentales para el diseño, optimización y control de los sistemas de secado, ya que permiten estimar tiempos de operación, mejorar la eficiencia energética y asegurar la calidad final del producto \parencite{fernando_effect_2025, ordonez-lozano_monitoring_2025}.

\subsubsection{Modelos empíricos y semiempíricos}

El modelamiento matemático del secado en materiales biológicos, como el café, es una herramienta esencial para predecir la evolución del contenido de humedad y optimizar las variables operativas del proceso \parencite{p_romero_leiton_mathematical_2025}. Estos modelos describen la relación entre humedad, tiempo y condiciones ambientales, facilitando el diseño y control de sistemas de secado \parencite{liu_optimization_2025}. En general, se clasifican en modelos empíricos y semiempíricos: los primeros se basan en regresiones de datos experimentales sin considerar explícitamente los mecanismos físicos, mientras que los segundos incorporan simplificaciones de leyes teóricas, como la segunda ley de Fick, logrando un balance entre precisión y simplicidad \parencite{khan_determination_2017, phitakwinai_thinlayer_2019}.

En el secado de café, se han validado diversos modelos para describir la reducción de la razón de humedad (MR). El modelo de Newton representa la aproximación más simple, mientras que el modelo de Page introduce un exponente que mejora el ajuste en procesos no lineales. Otros modelos, como Henderson y Pabis y el logarítmico, permiten describir con mayor precisión el comportamiento del material, especialmente cerca de la humedad de equilibrio \parencite{p_romero_leiton_mathematical_2025, phitakwinai_thinlayer_2019}. Los parámetros de estos modelos se estiman mediante regresión no lineal, reflejando la resistencia del material a la transferencia de humedad \parencite{cansee_effect_2026}.

La capacidad predictiva de los modelos se evalúa mediante indicadores estadísticos como el coeficiente de determinación ($R^2$), el error cuadrático medio (RMSE) y el chi-cuadrado ($\chi^2$) \parencite{abdissa_harmonizing_2023}. Un modelo adecuado presenta valores de $R^2$ cercanos a la unidad y errores mínimos, lo que garantiza una representación confiable del proceso \parencite{guemouni_convective_2022}. En este contexto, la selección de modelos robustos permite estimar tiempos de secado, optimizar el consumo energético y asegurar la calidad final del café \parencite{amaral_simulation_2018, p_romero_leiton_mathematical_2025}.

\subsubsection{Difusividad efectiva}

La difusividad efectiva ($D_{\text{eff}}$) es el parámetro que describe la velocidad de migración de la humedad desde el interior del material hacia su superficie durante el secado \parencite{hadibi_effect_2025, khan_determination_2017}. En materiales biológicos como el café, este fenómeno se modela a partir de la segunda ley de Fick, considerando un comportamiento global que integra distintos mecanismos de transporte, como difusión líquida, difusión de vapor, capilaridad y cambios estructurales del tejido \parencite{el-mesery_application_2025-1, cansee_effect_2026}. De este modo, $D_{\text{eff}}$ refleja la resistencia interna del material al movimiento de la humedad.

Este parámetro depende de factores térmicos, físicos y estructurales. La temperatura es el principal determinante, ya que incrementa la energía cinética del agua y favorece su movilidad, generalmente siguiendo una relación tipo Arrhenius \parencite{liu_optimization_2025, fernando_effect_2025}. Asimismo, la difusividad disminuye a medida que avanza el secado, debido a que el agua remanente se encuentra más fuertemente ligada a la matriz del grano. Factores como la porosidad, densidad y contracción del material también influyen en la resistencia al transporte de humedad \parencite{alves_physiological_2017, abdissa_harmonizing_2023}.

En el café, los valores de $D_{\text{eff}}$ varían según el estado del producto y las condiciones de secado, reportándose rangos del orden de $10^{-7}$ m$^2$/s en cerezas y entre $10^{-9}$ y $10^{-10}$ m$^2$/s en café pergamino \parencite{fernando_effect_2025, phitakwinai_thinlayer_2019}. La estimación precisa de este parámetro es fundamental para el modelamiento, diseño y optimización de secadores, ya que permite predecir la distribución de humedad, estimar tiempos de proceso y mejorar la eficiencia energética, evitando defectos como el endurecimiento superficial y preservando la calidad del grano \parencite{khan_determination_2017, cruz-ospina_improving_2026}.

\subsubsection{Ajuste estadístico y validación de modelos}

El ajuste estadístico y la validación de modelos matemáticos constituyen procesos fundamentales para calibrar las herramientas predictivas utilizadas en la cinética de secado de materiales biológicos como el café. El procedimiento general consiste en someter los datos experimentales de razón de humedad (MR) a un análisis de regresión no lineal, comúnmente mediante el método de mínimos cuadrados, con el fin de estimar los parámetros específicos que mejor describan el comportamiento del grano bajo condiciones operativas determinadas \parencite{fernando_effect_2025, guemouni_convective_2022, p_romero_leiton_mathematical_2025, phitakwinai_thinlayer_2019}. La calidad y exactitud de este ajuste se evalúan mediante indicadores de desempeño críticos: el coeficiente de determinación ($R^2$), el error cuadrático medio (RMSE) y el parámetro chi-cuadrado ($\chi^2$) \parencite{guemouni_convective_2022, abdissa_harmonizing_2023, cansee_effect_2026, meng_study_2023}. En términos de interpretación, un valor de $R^2$ cercano a la unidad indica que el modelo captura con éxito la mayor parte de la variabilidad de los datos, mientras que valores mínimos de RMSE y $\chi^2$ (tendientes a cero) confirman una desviación reducida entre los valores predichos y los observados, garantizando la superioridad del modelo para representar el fenómeno físico del secado \parencite{phitakwinai_thinlayer_2019, fernando_effect_2025, guemouni_convective_2022, p_romero_leiton_mathematical_2025}.

La validación de estos modelos mediante la comparación con datos independientes o experimentales es un paso indispensable para asegurar su robustez y aplicabilidad en contextos reales \parencite{amaral_simulation_2018, fauzi_experimental_2025, meng_study_2023}. Este proceso permite verificar si las predicciones matemáticas, como las derivadas de los modelos de Midilli o Page, son consistentes con la dinámica de transferencia de masa y calor observada en el equipo de secado, incluso frente a fluctuaciones en la temperatura y la humedad relativa \parencite{amaral_simulation_2018, p_romero_leiton_mathematical_2025, phitakwinai_thinlayer_2019}. En la ingeniería de procesos, una validación exitosa otorga la confiabilidad necesaria para el diseño y control de los sistemas industriales, permitiendo optimizar variables críticas sin recurrir a pruebas físicas exhaustivas y costosas \parencite{fauzi_experimental_2025, leonardo_barcelos_mateus_modeling_2021}. 


\subsection{Impacto del secado en la calidad física y sensorial}

\subsubsection{Actividad de agua y estabilidad}

La actividad de agua ($a_w$) es un parámetro termodinámico que mide la fracción de agua disponible en el grano para reacciones químicas y crecimiento microbiano, diferenciándose del contenido de humedad, que representa la cantidad total de agua \parencite{abreu_influence_2025, penuela_martinez_influence_2023}. Ambos parámetros están correlacionados, de modo que un mayor contenido de humedad implica mayor $a_w$ y, por tanto, mayor susceptibilidad al deterioro \parencite{abreu_influence_2025}. Por esta razón, la $a_w$ es un indicador clave de estabilidad en productos alimentarios \parencite{largoavila_changes_2023}.

La estabilidad microbiológica y química del café depende de mantener la $a_w$ en niveles seguros. Valores superiores a 0.5 favorecen el crecimiento de hongos como \textit{Aspergillus}, \textit{Penicillium} y \textit{Fusarium}, así como la posible formación de micotoxinas \parencite{abreu_influence_2025}. Además, una alta $a_w$ promueve reacciones de deterioro como oxidación, hidrólisis y rancidez lipídica. En general, se recomienda mantener la $a_w$ por debajo de 0.6 para asegurar la estabilidad del producto durante almacenamiento \parencite{penuela_martinez_influence_2023}.

En el secado de café, esto se logra al reducir la humedad hasta valores cercanos al 10--12\%, lo que corresponde a una $a_w$ aproximada de 0.59 \parencite{guevara-sanchez_effect_2019, huaman-murillo_physical_2024}. Un control adecuado del proceso evita tanto el subsecado, que incrementa el riesgo microbiológico, como el sobresecado, que puede afectar la calidad sensorial y la integridad del grano \parencite{abdissa_harmonizing_2023, sandeep_effect_2021}. Por ello, el monitoreo de la $a_w$ es fundamental para garantizar la estabilidad, inocuidad y calidad del café de especialidad.

\subsubsection{Efectos térmicos en compuestos precursores aromáticos}

La calidad sensorial del café depende de la composición química del grano verde, constituida por compuestos precursores como carbohidratos (principalmente sacarosa), lípidos, proteínas, aminoácidos, trigonelina y ácidos clorogénicos. Durante la tostión, estos compuestos participan en reacciones como Maillard y caramelización, generando numerosos compuestos volátiles responsables del aroma, cuerpo y sabor de la bebida \parencite{eshetu_effect_2022, tieghi_effects_2024}. En este contexto, los lípidos contribuyen a la retención de volátiles, mientras que la trigonelina aporta notas características tras su degradación térmica \parencite{largoavila_changes_2023}.

El secado térmico influye directamente en la concentración y estabilidad de estos precursores antes de la tostión. Condiciones controladas y tiempos de secado adecuados favorecen la conservación de compuestos clave como la sacarosa y los ácidos grasos, mientras que procesos prolongados pueden inducir su degradación debido a la actividad metabólica residual del grano, afectando negativamente el potencial sensorial \parencite{largoavila_changes_2023, borem_quality_2018}.

El control de la temperatura es especialmente crítico, ya que valores superiores a 40--50~°C pueden generar daños estructurales en las membranas celulares, promoviendo la oxidación de lípidos y la formación de defectos sensoriales (\textit{off-flavors}) \parencite{borem_quality_2018, oliveira_quality_2018}. Asimismo, el exceso de calor acelera la degradación de ácidos clorogénicos y compuestos nitrogenados, reduciendo la complejidad aromática y aumentando el riesgo de sabores amargos o quemados \parencite{maryana_performance_2025}. Por ello, el secado a temperaturas moderadas (preferiblemente menores a 40~°C) es fundamental para preservar la calidad sensorial y cumplir con los estándares de cafés de especialidad \parencite{abreu_influence_2025, de_almeida_rios_origin_2021}.

\subsubsection{Requerimientos de precisión en cafés de especialidad}

 En el mercado de especialidad, la ausencia total de defectos primarios y la mínima presencia de secundarios son estándares innegociables que dependen directamente de un secado homogéneo y controlado \parencite{escuela_superior_politecnica_agropecuaria_de_manabi_manuel_felix_lopez_carrera_de_ingenieria_agricola_postharvest_2023, freitas_santos_efect_2020, huaman-murillo_physical_2024, largo-avila_influence_2023}.

El cumplimiento de estos requerimientos de precisión tiene una repercusión directa en el valor comercial y la competitividad del caficultor en mercados internacionales \parencite{barbosa_sensory_2019, carmo_quality_2020, freitas_santos_efect_2020}. La implementación de tecnologías de secado que optimizan el flujo de aire y la transferencia de calor reduce la incidencia de granos quebradizos o descoloridos, incrementando el rendimiento de café exportable y la puntuación final en taza \parencite{abreu_influence_2025, konopatzki_coffee_2022, konopatzki_price_2019}. Una mejora en la calificación sensorial no solo posiciona al café en categorías superiores (como Q Premium o Exceptional), sino que también puede representar un incremento significativo en el precio de venta, reportándose valorizaciones superiores al 12\% en comparación con sistemas de secado tradicionales menos precisos \parencite{firdissa_coffee_2022, abreu_influence_2025, huaman-murillo_physical_2024, konopatzki_price_2019}. En consecuencia, la gestión científica del secado se consolida como el factor diferenciador que transforma una materia prima en un producto de alto valor agregado y calidad sostenible \parencite{carmo_quality_2020, freitas_santos_efect_2020, largo-avila_influence_2023}.

\subsection{Tecnologías de secado aplicadas al café}

\subsubsection{Secado solar}

El secado solar de café se define como un proceso de deshidratación térmica fundamentado en el aprovechamiento de la radiación solar como fuente primaria de energía para reducir el contenido de humedad del grano hasta niveles que garanticen su conservación y estabilidad \parencite{firdissa_coffee_2022, guevara-sanchez_effect_2019, kebede_performance_2025}. Este fenómeno implica una transferencia simultánea de calor y masa, donde el agua interna se difunde hacia la superficie para ser evaporada por el aire circundante, proceso que puede potenciarse mediante el efecto invernadero en estructuras cubiertas para elevar la temperatura del aire \parencite{firdissa_coffee_2022, siagian_field_2017}. Entre los métodos más extendidos se encuentran el secado tradicional en patios de cemento o mantas directamente sobre el suelo, el uso de camas africanas o parihuelas elevadas que facilitan el flujo de aire por convección natural, y estructuras más tecnificadas como las marquesinas o secadores solares tipo túnel e invernadero \parencite{firdissa_coffee_2022, huaman-murillo_physical_2024}. Esta tecnología es valorada académicamente por su bajo costo de instalación, estructura simple y sostenibilidad, al emplear una energía limpia que no contamina el medio ambiente, lo que la hace altamente accesible para pequeños productores en diversas condiciones geográficas \parencite{guevara-sanchez_effect_2019, kebede_performance_2025}.

No obstante, desde una perspectiva técnica y crítica, el secado solar presenta limitaciones significativas que condicionan su eficiencia frente a sistemas mecánicos. La principal desventaja es su absoluta dependencia climática, lo que vuelve el proceso vulnerable a la nubosidad y precipitaciones imprevistas que pueden provocar la rehumectación del grano y generar pérdidas económicas considerables en regiones tropicales \parencite{guevara-sanchez_effect_2019, huaman-murillo_physical_2024, kebede_performance_2025}. Estas condiciones climáticas fluctuantes derivan en tiempos prolongados de secado, que en métodos tradicionales pueden extenderse por ocho días o más, incrementando el riesgo de deterioro microbiológico si la humedad permanece por encima del 12.5\% \parencite{guevara-sanchez_effect_2019, firdissa_coffee_2022, huaman-murillo_physical_2024}. Asimismo, la dificultad de control preciso sobre la temperatura y la humedad relativa dentro de la masa de café suele generar un secado heterogéneo y desigual. Por último, la exposición en sistemas abiertos conlleva un elevado riesgo de contaminación por polvo, microorganismos, insectos y animales, lo que compromete la inocuidad del grano \parencite{huaman-murillo_physical_2024, kebede_performance_2025}.

Estas deficiencias técnicas impactan directamente en la calidad final del café, correlacionándose con la aparición de defectos físicos (granos agrios, negros o con moho) y la pérdida de atributos sensoriales en taza \parencite{firdissa_coffee_2022, huaman-murillo_physical_2024}. La falta de uniformidad térmica produce lotes con heterogeneidad en el secado, lo que puede derivar en una actividad bioquímica indeseada que genera sabores desagradables durante el almacenamiento \parencite{guevara-sanchez_effect_2019, huaman-murillo_physical_2024}. De este modo, aunque el secado solar es una solución accesible y sostenible, su implementación requiere una gestión rigurosa para evitar que la falta de control ambiental degrade el perfil de los cafés especiales, los cuales demandan una consistencia física y sensorial superior para alcanzar precios diferenciados en el mercado mundial \parencite{guevara-sanchez_effect_2019, huaman-murillo_physical_2024}.

\subsubsection{Secado mecánico convencional}

El secado mecánico convencional se fundamenta en un proceso convectivo que emplea fuentes de calor artificial y ventilación forzada para la remoción de humedad mediante la transferencia simultánea de masa y energía \parencite{abreu_influence_2025, coradi_development_2021, largo-avila_influence_2023}. Este método utiliza equipos especializados tales como secadores de lecho o capa fija (estáticos), secadores rotatorios (comúnmente denominados guardiolas), silos-secadores y sistemas de flujo continuo, los cuales permiten un contacto directo del aire caliente con la masa del grano \parencite{alves_physiological_2017, coradi_development_2021, leonardo_barcelos_mateus_modeling_2021, prada_efectividad_2019, sandeep_effect_2021}. Frente al secado solar tradicional, la tecnificación mecánica ofrece ventajas operativas críticas: una reducción sustancial del tiempo de proceso —logrando alcanzar el contenido de humedad óptimo (10--12\%) en lapsos de 24 a 48 horas frente a las semanas que requiere el método natural—, una nula dependencia de las condiciones climáticas adversas y una mayor capacidad de procesamiento necesaria para gestionar los altos volúmenes de café provenientes de cosechas mecanizadas \parencite{abreu_influence_2025, alves_physiological_2017, largo-avila_influence_2023, sandeep_effect_2021}.

A pesar de su eficiencia logística, el secado mecánico convencional presenta limitaciones técnicas y económicas significativas. Es la etapa poscosecha de mayor demanda energética, representando a menudo más del 60\% del consumo total de energía en la producción, lo que conlleva altos costos operativos asociados al uso de combustibles como GLP o biomasa \parencite{coradi_development_2021}. Críticamente, muchos sistemas comerciales operan con una baja eficiencia térmica y presentan dificultades para el control preciso de variables ambientales, especialmente la humedad relativa del aire de secado, que influye directamente en la tasa de evaporación \parencite{coradi_development_2021, andrade_mathematical_2019, fernando_effect_2025}. Esta falta de control riguroso incrementa el riesgo de sobrecalentamiento del grano; la práctica común de elevar las temperaturas por encima de los límites recomendados para acelerar el proceso suele resultar en daños irreversibles en la estructura biológica del café \parencite{alves_physiological_2017, borem_quality_2018, coradi_development_2021}.

Las deficiencias en el manejo térmico tienen una relación directa y negativa con la calidad sensorial y física del producto final. El uso de temperaturas en la masa del grano superiores a 40--45~°C provoca el colapso y la desorganización de las membranas celulares del endospermo, lo que desencadena procesos de lixiviación de solutos y oxidación de lípidos \parencite{borem_quality_2018, coradi_development_2021, oliveira_quality_2018, penuela_martinez_influence_2023, sandeep_effect_2021}. Estos daños térmicos se manifiestan en defectos físicos como el blanqueamiento o el oscurecimiento del grano, así como en la pérdida de compuestos volátiles esenciales para el aroma y el sabor \parencite{abreu_influence_2025, coradi_development_2021, parra-coronado_preliminary_2020, penuela_martinez_influence_2023}. En consecuencia, un secado mecánico mal ejecutado no solo degrada los atributos organolépticos (reduciendo cuerpo, fragancia y dulzor), sino que también compromete la estabilidad química del café durante el almacenamiento, propiciando la aparición de notas fermentadas o rancias que devalúan comercialmente el lote \parencite{borem_quality_2018, sandeep_effect_2021, alves_physiological_2017, coradi_development_2021, oliveira_quality_2018, parra-coronado_preliminary_2020}.

\subsubsection{Secado asistido por bomba de calor}

El sistema de secado asistido por bomba de calor (SBC) se define como la integración de dos subsistemas termodinámicos: un secador convectivo y un ciclo de refrigeración por compresión de vapor \parencite{gomez-daza_innovacion_2023}. Su principio de funcionamiento se basa en el ciclo de Carnot inverso, integrando cuatro componentes esenciales vinculados por una línea de refrigerante: un compresor, un condensador, una válvula de expansión y un evaporador \parencite{fauzi_influence_2024, salehi_recent_2021}. En este proceso, el aire húmedo que sale de la cámara de secado pasa por el evaporador, el cual actúa como un deshumidificador; allí, el aire se enfría sensiblemente hasta alcanzar su punto de rocío, permitiendo que el vapor de agua se condense y sea retirado del sistema \parencite{salehi_recent_2021, dzaky_activation_2022}. El calor latente de vaporización liberado durante esta condensación es absorbido por el refrigerante y transferido mediante el compresor hacia el condensador, donde se utiliza para calentar sensiblemente el aire ya seco antes de que éste reingrese a la cámara de secado \parencite{salehi_recent_2021, yahya_development_2023}. Este diseño en circuito cerrado permite una recuperación integral del calor latente y sensible, lo que constituye la base de su alta eficiencia energética en comparación con sistemas de ventilación abierta \parencite{salehi_recent_2021}.

Una de las ventajas más relevantes del SBC es su capacidad para realizar un control preciso e independiente del estado psicrométrico del aire, regulando simultáneamente la temperatura y la humedad relativa \parencite{salehi_recent_2021, gomez-daza_innovacion_2023}. A diferencia del secado solar, que depende de las fluctuaciones climáticas y conlleva riesgos de contaminación o fermentación por tiempos prolongados, o del secado mecánico convencional, que a menudo requiere altas temperaturas y un elevado consumo energético, el secador asistido por bomba de calor permite operar en un amplio rango de condiciones ambientales con un control fino del proceso \parencite{salehi_recent_2021, cheng_evaluation_2019, dong_simulation_2022}. Este sistema garantiza una mayor uniformidad en el secado al mantener un potencial de secado constante, definido por la presión de vapor en la superficie del producto y la presión parcial del vapor de agua en el aire \parencite{gomez-daza_innovacion_2023}. Además, el uso de válvulas de expansión electrónicas y compresores de velocidad variable permite optimizar el coeficiente de desempeño (COP) y la velocidad específica de extracción de humedad (SMER), logrando ahorros energéticos de entre el 60\% y el 80\% respecto a los secadores convencionales que operan a la misma temperatura \parencite{salazar-hincapie_experimental_2020, salehi_recent_2021}.

Desde la perspectiva de la calidad del producto, el secado asistido por bomba de calor es especialmente adecuado para materiales termo-sensibles como el café, que pueden sufrir daños estructurales, térmicos o pérdida de aroma si se exponen a condiciones severas \parencite{salehi_recent_2021, fauzi_influence_2024}. El control preciso de las variables de operación minimiza el riesgo de endurecimiento superficial y oscurecimiento no deseado, preservando mejor los compuestos bioactivos, como los ácidos clorogénicos y la actividad antioxidante, en comparación con métodos de secado menos controlados que degradan estos elementos debido a tiempos de exposición prolongados \parencite{cheng_evaluation_2019}. Asimismo, se ha reportado que los perfiles volátiles y sensoriales del café procesado mediante tecnologías de secado controlado pueden mantener una alta retención de los aromas frescos, evitando defectos como el moho o la rancidez \parencite{salehi_recent_2021}. En consecuencia, la integración de la bomba de calor representa una alternativa tecnológica avanzada que asegura la integridad química y sensorial del grano de café, optimiza el rendimiento energético y reduce el impacto ambiental de la operación poscosecha \parencite{salehi_recent_2021, yahya_development_2023}.

\subsection{Evaluación energética y eficiencia en sistemas de secado}

\subsubsection{Balance energético}

El balance energético en sistemas de secado de café, fundamentado en la primera ley 
de la termodinámica, establece que la energía total ingresada —de tipo térmico 
(calentadores eléctricos, radiación solar o combustión de biomasa y gas) y eléctrico 
(ventiladores y sistemas de control)— debe equipararse a la energía de salida más 
los cambios en el almacenamiento interno 
\parencite{leonardo_barcelos_mateus_modeling_2021, rabello_modeling_2021, 
suherman_energy_2020, abbaspour-gilandeh_effect_2021, arifin_controlling_2024, 
yuwana_free_2025}. La energía suministrada se distribuye entre el calentamiento del 
aire (calor sensible), el calentamiento del material y, de forma predominante, la 
evaporación del agua —que en sistemas optimizados consume más del 69\% de la energía 
total—; no obstante, pérdidas significativas ocurren a través de las paredes de la 
cámara, ineficiencias mecánicas y, especialmente, por el calor transportado en el 
aire de salida expulsado sin aprovechamiento 
\parencite{burdo_energy-efficient_2025, santana_economic_2020, 
abbaspour-gilandeh_effect_2021, li_study_2021}.

La eficiencia térmica, definida como la relación entre la energía utilizada para 
evaporar la humedad y la energía total consumida, es reducida en sistemas convectivos 
convencionales debido al desperdicio de calor en el flujo de salida 
\parencite{abbaspour-gilandeh_effect_2021, gomez-daza_innovacion_2023, li_study_2021, 
burdo_energy-efficient_2025}. En contraste, el secado asistido por bomba de calor 
opera en ciclo cerrado recuperando el calor latente del aire húmedo, minimizando 
pérdidas y reduciendo el consumo específico de energía; este enfoque no solo optimiza 
el rendimiento energético sino que favorece un control más riguroso de las variables 
térmicas, preservando las propiedades sensoriales y la calidad final del café 
\parencite{burdo_energy-efficient_2025, gomez-daza_innovacion_2023, pramono_low_2022, 
rabello_modeling_2021}.

\subsubsection{Consumo energético específico (SEC)}

El consumo energético específico (SEC) cuantifica la energía total requerida para 
eliminar un kilogramo de agua durante el secado, constituyendo un indicador clave 
de eficiencia térmica que permite comparar tecnologías y condiciones operativas 
independientemente de la escala \parencite{burdo_energy-efficient_2025, 
el-mesery_application_2025-1, liu_optimization_2025, gomez-daza_innovacion_2023, 
institute_of_food_and_biotechnology_can_tho_university_vietnam_impact_2022}. En el 
procesamiento de café, donde el secado representa una de las etapas más intensivas 
en energía, un SEC menor indica mayor eficiencia y menores costos de producción 
\parencite{cansee_effect_2026, cruz-ospina_improving_2026, 
leonardo_barcelos_mateus_modeling_2021, chojnacka_improvements_2021, 
kaveh_comparative_2023}. Los principales factores que lo afectan son la temperatura 
del aire —cuyo incremento reduce el SEC al acelerar la transferencia de masa—, la 
velocidad del aire —con efecto dual, ya que flujos excesivos generan pérdidas 
convectivas— y la humedad relativa —donde valores bajos intensifican la fuerza 
impulsora de evaporación—, siendo el periodo de velocidad decreciente la etapa más 
limitante por la difusión interna del agua 
\parencite{institute_of_food_and_biotechnology_can_tho_university_vietnam_impact_2022, 
abbaspour-gilandeh_effect_2021, kaveh_comparative_2023, el-mesery_application_2025-1, 
liu_optimization_2025, cruz-ospina_improving_2026, gomez-daza_innovacion_2023, 
li_study_2021, yu_current_2025}.

El análisis comparativo del SEC evidencia diferencias significativas entre 
tecnologías: el secado solar presenta bajo consumo energético directo pero depende 
de condiciones ambientales; los secadores mecánicos convectivos registran valores 
más altos por pérdidas de calor sensible en el aire de salida; y el secado asistido 
por bomba de calor alcanza los valores más bajos al recuperar el calor latente de 
vaporización \parencite{cansee_effect_2026, burdo_energy-efficient_2025, 
ibrahim_design_2025, gomez-daza_innovacion_2023}. La reducción del SEC se vincula 
directamente con la sostenibilidad del proceso, al traducirse en menores costos 
operativos y en una disminución de emisiones de gases de efecto invernadero 
\parencite{cansee_effect_2026, kaveh_comparative_2023}.

\subsubsection{Tasa específica de extracción de humedad (SMER)}

El SMER se define como la masa de agua removida por unidad de energía consumida, 
constituyendo el recíproco del SEC; maximizar el SMER equivale directamente a 
reducir el SEC, facilitando la comparación técnica entre tecnologías de secado 
independientemente de su escala \parencite{burdo_energy-efficient_2025, 
gomez-daza_innovacion_2023, hendrarsakti_thermal_2019, 
leonardo_barcelos_mateus_modeling_2021}. Su magnitud es proporcional a la 
temperatura, la humedad inicial del material y la carga del secador; un aspecto 
técnico distintivo es que en sistemas de ciclo cerrado el SMER tiende a aumentar 
con la humedad relativa (HR) del aire, en contraste con la velocidad de secado, 
favorecida por condiciones de baja HR y alta velocidad del aire, lo que exige 
balancear estos factores en el diseño de estrategias de control para maximizar la 
extracción de agua sin comprometer la calidad sensorial del grano 
\parencite{gomez-daza_innovacion_2023}.

Los sistemas de secado asistido por bomba de calor (SBC) alcanzan los valores más 
altos de SMER al recuperar el calor sensible y latente del aire húmedo para 
reintegrarlo al proceso, superando las ineficiencias de los secadores convectivos 
convencionales que desperdician calor en el flujo de salida 
\parencite{burdo_energy-efficient_2025, abbaspour-gilandeh_effect_2021, 
chojnacka_improvements_2021}. Estudios experimentales en materiales agroalimentarios 
reportan valores de SMER para SBC entre 0.156 y 9.25~kg/kWh, superando 
notablemente a sistemas solares pasivos y mecánicos de combustión, lo que posiciona 
esta tecnología como una alternativa clave para el desarrollo sostenible de la 
cadena productiva del café \parencite{bobbo_technological_2024, gomez-daza_innovacion_2023, 
burdo_energy-efficient_2025}.

\subsubsection{Indicadores de eficiencia energética}

La eficiencia energética en el secado de café se evalúa mediante tres indicadores 
complementarios: el balance energético —que distribuye la energía ingresada entre 
evaporación, calentamiento del aire y del grano, y pérdidas al ambiente—, el SEC 
—que cuantifica la energía requerida para eliminar un kilogramo de agua— y el SMER 
—su recíproco, que indica la masa de agua removida por unidad de energía consumida— 
\parencite{burdo_energy-efficient_2025, gomez-daza_innovacion_2023, 
geng_thermodynamically_2023, abbaspour-gilandeh_effect_2021, liu_optimization_2025}. 
La relación inversa entre SEC y SMER permite que la optimización de uno mejore 
directamente el otro, facilitando la comparación objetiva entre tecnologías y escalas 
de operación \parencite{gomez-daza_innovacion_2023, burdo_energy-efficient_2025}. 
Dado que el secado es la etapa más intensiva en energía del procesamiento de café, 
un SEC elevado incrementa los costos de producción y la dependencia de combustibles 
fósiles, mientras que su reducción disminuye la huella de carbono y promueve una 
producción más limpia en concordancia con estándares internacionales de sostenibilidad 
\parencite{cruz-ospina_improving_2026, el-mesery_application_2025-1, 
kaveh_comparative_2023, institute_of_food_and_biotechnology_can_tho_university_vietnam_impact_2022, 
cansee_effect_2026, chojnacka_improvements_2021}.

El secado asistido por bomba de calor representa la tecnología de mayor desempeño 
en estos indicadores, al operar en ciclo cerrado recuperando el calor sensible y 
latente del aire húmedo —a diferencia de los sistemas convectivos convencionales 
que lo desperdician en el flujo de salida— \parencite{bobbo_technological_2024, 
chojnacka_improvements_2021}. Esta recuperación permite alcanzar valores de SMER 
superiores a 9.0~kg/kWh y reducir la demanda energética total hasta en un 80\%, 
al tiempo que el control preciso de temperatura y humedad preserva las propiedades 
sensoriales, la estabilidad fisiológica y la calidad final del café pergamino 
\parencite{burdo_energy-efficient_2025, bobbo_technological_2024, 
el-mesery_application_2025, gomez-daza_innovacion_2023}.

\subsection{Escalamiento y validación tecnológica}

\subsubsection{Principios de escalamiento térmico}

El escalamiento en los procesos de secado representa la transición técnica desde experimentos controlados a nivel de laboratorio o planta piloto hacia la producción a escala industrial \parencite{szpicer_advances_2025}. Este proceso no es una expansión lineal, ya que al aumentar el tamaño del sistema, la complejidad geométrica se incrementa y los regímenes de flujo pueden cambiar, por ejemplo, de laminar a turbulento, lo que provoca que las dinámicas del sistema se vuelvan significativamente no lineales. Los principales desafíos térmicos del escalamiento radican en la sensibilidad de los procesos a factores dependientes de la escala, tales como los cambios en los mecanismos de transferencia de calor y masa, la generación de calor y la dificultad para mantener una distribución uniforme de temperatura y humedad en grandes volúmenes \parencite{szpicer_advances_2025}. En sistemas industriales, un diseño deficiente puede generar zonas de estancamiento o recirculación de aire, lo que incrementa el tiempo de procesamiento y la demanda energética global \parencite{szpicer_advances_2025, yu_current_2025}.

Para asegurar que las propiedades del producto se mantengan durante el escalamiento, es fundamental el diseño basado en principios de similitud geométrica, térmica y dinámica \parencite{szpicer_advances_2025}. Aunque tradicionalmente se emplean leyes de escalamiento basadas en números adimensionales, su aplicación puede presentar limitaciones en sistemas de matrices alimentarias complejas y multifásicas como el café \parencite{szpicer_advances_2025}. Variables operacionales como el flujo de aire, el tamaño de los componentes del sistema y la carga de material afectan drásticamente el comportamiento térmico; por ejemplo, aberturas adecuadamente dimensionadas en los conductos facilitan un flujo de aire más equilibrado, reduciendo las pérdidas de energía asociadas a irregularidades en la velocidad \parencite{du_design_2025}. Asimismo, la densidad de carga del material es crítica, ya que un aumento excesivo en la cantidad de producto puede alterar el campo de flujo, incrementando la resistencia a la remoción de humedad y afectando negativamente la transferencia de masa convectiva \parencite{onwude_scaling-up_2021, cansee_effect_2026}.

Finalmente, el éxito del escalamiento térmico está intrínsecamente relacionado con la eficiencia energética y la preservación de la calidad del café \parencite{cruz-ospina_improving_2026, szpicer_advances_2025}. Un escalamiento optimizado busca minimizar el secado no uniforme, el cual puede causar sobresecado —que degrada compuestos de interés y aumenta la fragilidad del grano— o subsecado, que favorece el crecimiento microbiano durante el almacenamiento \parencite{yu_current_2025}. La mejora en la retención de calor mediante el diseño optimizado de la cámara o el uso de materiales aislantes permite reducir los tiempos de proceso, incrementar la capacidad de producción y disminuir los costos operativos unitarios \parencite{cruz-ospina_improving_2026}. En última instancia, un enfoque de escalamiento que integre el modelado físico y la validación experimental permite desarrollar soluciones que equilibran productividad, eficiencia energética y estabilidad fisicoquímica del grano de café \parencite{du_design_2025, yu_current_2025}.

\subsubsection{Variables críticas de diseño}

El diseño de sistemas de secado para café se basa en la interacción de variables termodinámicas y estructurales que gobiernan la transferencia de calor y masa. La temperatura del aire y la humedad relativa constituyen las principales fuerzas impulsoras del proceso, influyendo directamente en la velocidad de secado y en la difusividad efectiva \parencite{geng_thermodynamically_2023, yu_current_2025}. Asimismo, la velocidad y distribución del aire determinan la eficiencia del transporte convectivo; mientras que condiciones adecuadas reducen la resistencia de la capa límite, flujos mal distribuidos generan pérdidas de energía y secado no uniforme \parencite{cruz-ospina_improving_2026, liu_optimization_2025}.

La geometría del secador y la carga de material (espesor del lecho) condicionan la uniformidad del proceso y el desempeño energético. Diseños adecuados de plenum y conductos permiten un flujo homogéneo y evitan zonas de estancamiento o puntos calientes. Por su parte, una mayor carga incrementa la resistencia al flujo, pero puede mejorar el aprovechamiento energético reflejado en indicadores como el SMER, lo que evidencia la necesidad de un balance óptimo entre capacidad y eficiencia \parencite{du_design_2025, cansee_effect_2026}.

Finalmente, la integración de estas variables en sistemas de control es clave para la estabilidad del proceso. El uso de controladores PID, sensores en tiempo real y tecnologías como el secado asistido por bomba de calor permite regular de manera precisa el estado psicrométrico del aire, mejorando la eficiencia y la calidad del producto. Adicionalmente, herramientas como la dinámica de fluidos computacional (CFD) facilitan el diseño y optimización de los equipos al predecir el comportamiento térmico y del flujo \parencite{arifin_controlling_2024, rabello_modeling_2021}.

\subsubsection{Validación experimental}

La validación experimental en los sistemas de secado de café constituye un proceso crítico para garantizar la precisión de los modelos matemáticos que describen la cinética de remoción de humedad \parencite{briceno-martinez_propuesta_2020, cruz-ospina_improving_2026}. Este procedimiento implica la comparación sistemática de los datos obtenidos en campo o planta piloto con las predicciones teóricas basadas en modelos empíricos o semiempíricos, tales como los de Page, Midilli o Newton, evaluando su ajuste mediante indicadores estadísticos como el coeficiente de determinación ($R^2$), el error cuadrático medio (RMSE) y el parámetro chi-cuadrado ($\chi^2$) \parencite{cansee_effect_2026, guemouni_convective_2022, institute_of_food_and_biotechnology_can_tho_university_vietnam_impact_2022, pramono_low_2022, suherman_energy_2020}. Las variables fundamentales medidas durante estos ensayos incluyen el contenido de humedad del grano —el cual debe reducirse habitualmente desde niveles de cosecha cercanos al 55\% hasta un rango de estabilidad biológica de 11--12\% b.h.—, la temperatura del aire y del material, la humedad relativa del ambiente y el tiempo total de secado \parencite{cruz-ospina_improving_2026, hendrarsakti_thermal_2019, leonardo_barcelos_mateus_modeling_2021, pramono_low_2022}. Esta correlación entre datos experimentales y modelos es esencial para predecir el comportamiento del sistema bajo diversas condiciones operativas, asegurando que la calidad final y la estabilidad fisiológica del café se mantengan dentro de los estándares comerciales \parencite{leonardo_barcelos_mateus_modeling_2021, cruz-ospina_improving_2026, rabello_modeling_2021, suherman_energy_2020}.

La robustez de la validación experimental depende estrictamente de la reproducibilidad y repetibilidad de los ensayos, los cuales suelen realizarse mediante múltiples réplicas para minimizar la incertidumbre y los errores aleatorios inherentes a los procesos biológicos \parencite{cruz-ospina_improving_2026, abbaspour-gilandeh_effect_2021, cansee_effect_2026, du_design_2025, liu_optimization_2025}. Para lograr este nivel de precisión, es indispensable el uso de herramientas tecnológicas como sensores de temperatura y humedad, sistemas de adquisición de datos y monitoreo en tiempo real \parencite{ambarita_design_2020, el-mesery_application_2025-1, geng_thermodynamically_2023, pramono_low_2022, yu_current_2025}. Tecnologías emergentes basadas en el Internet de las Cosas (IoT) permiten registrar de manera continua variables críticas mediante termocuplas, higrómetros digitales y sistemas de pesaje del producto, facilitando un control dinámico del proceso a través de estrategias automatizadas \parencite{pramono_low_2022, arifin_controlling_2024, cruz-ospina_improving_2026, yu_current_2025}. 

Finalmente, la validación experimental no solo corrobora la eficiencia técnica del diseño, sino que permite calcular con precisión indicadores de desempeño como el consumo energético específico (SEC) y la tasa específica de extracción de humedad (SMER) \parencite{cruz-ospina_improving_2026, burdo_energy-efficient_2025, cansee_effect_2026, gomez-daza_innovacion_2023}. Estos indicadores permiten comparar objetivamente diferentes tecnologías de secado, asegurando que la optimización del sistema se traduzca en una reducción efectiva del consumo energético y en una mejora tangible de la sostenibilidad del proceso industrial \parencite{cansee_effect_2026, cruz-ospina_improving_2026, gomez-daza_innovacion_2023, guemouni_convective_2022}.